\documentclass[a4paper]{article}
\usepackage[pdftex]{graphicx}
\usepackage[utf8]{inputenc}
\usepackage{enumerate}
\usepackage{icomma}
\usepackage{amssymb}
\usepackage{tikz}
\usepackage{href-ul}
\hypersetup{
	colorlinks=true,
	linkcolor=blue,
	urlcolor=blue}
\usepackage{geometry}
\geometry{a4paper, top=15mm, left=15mm, right=15mm, bottom=15mm,
	headsep=10mm, footskip=12mm}

\begin{document}
\begin{tikzpicture}(10,3)
	\draw[ultra thick](2,0) --(10,0) -- (10,1.3) --(2,1.3) -- (2,0);
	\draw[fill=black](2,0)-- (0,.6) -- (2,.6) -- (2,0);
	\node at (6,.6) {\large Lösungsblatt von \href{https://www.okuyakl.de}{www.okuyakl.de}};
\end{tikzpicture}
\vspace{0.5 cm}

\noindent {\bf Aufgabe 1. a)} \\
Erste  Ableitung mit Kettenregel:
$$f'(x)=-2\ln{x}\cdot{1 \over x}$$
Zweite Ableitung mit Produktregel: 
$$
\begin{array}{rcl}
f''(x)&=&{2 \over x^2} \cdot \ln{x} - {2 \over x} \cdot {1\over x}\\
\ &=& {2 \over x^2}  \cdot \left(\ln{x}-1 \right)\stackrel{!}{=}0\\
\ln{x_0}&=&1 \\
x_0&=&e
\end{array}
$$

Tangente: $$y_T=m \cdot x+t; \qquad x=e; \qquad m=f'(e)= -{2 \over e}; \qquad y_T= f(e)=0 $$
$$ 0=-{2 \over e} \cdot e + t; \quad \Rightarrow t=2$$
$$ \Rightarrow y= -{2 \over e}\cdot x +2 $$

\noindent {\bf  Aufgabe 1. b)} \\
$F$ mit Produkt- und Kettenregel ableiten:
$$
\begin{array}{rcl}
F'(x)&=&-(\ln{x}-1)^2 - x \cdot 2 \cdot (\ln{x}-1) \cdot {1 \over x}\\
\ &=& -((\ln{x})^2 - 2 \ln{x}+1) - 2 \ln{x} + 2 \\
\ &=& 1-(\ln{x})^2
\end{array}
$$

\noindent {\bf  Aufgabe 1. c)}\\
Nullstellen: 
$$
\begin{array}{rcl}
1-(\ln{x})^2 &=& 0 \\
\ln{x} &=& \pm 1 \\
x_1&=&e \\
x_2&=&e^{-1}
\end{array}
$$
Damit und mit der gegebenen Stammfunktion berechnen wir die Fläche:
$$A=\int\limits_{e^{-1}}^e f(x) ~{\rm d} x =\left[ -x(\ln x -1)^2\right]_{e^{-1}}^e = 0 -(-e^{-1}(-2)^2) = {4\over e}$$

\noindent {\bf   Aufgabe 1. d)}\\
Jede Integralfunktion hat eine Nullstelle, die dort liegt, wo die obere Grenze, hier die Variable $x$ mit der unteren Grenze, hier die 0, zusammenfällt, also bei $x=0$. Dann ist $I(0)=0$, hier ist
auch $ F(0)=0$, also ist $F(x)=I(x)$

\noindent {\bf  Aufgabe 1. e)}\\
Laut Merkhilfe ist der Grenzwert $$\lim \limits_{x \to 0^+} x \cdot ln{x} =0 $$
und damit auch:
$$\lim \limits_{x \to 0^+} -x \cdot ( ln{x}-1)^2 =0 $$
Das bedeutet, dass das Flächenstück, das von der Funktion und den beiden Achsen im IV. Quadranten begrenzt wird, endlich ist. 

\noindent {\bf  Aufgabe 2. a)}\\
Es gilt: $I'(x)=f(x)$, also setzen wir $f(x)=0$:
$$
\begin{array}{rcl}
(e^x-2)^2 &=& 0\\
(e^x-2)&=& 0 \\
e^x &=& 2 \\
x&=& \ln{2}
\end{array}
$$
Monotonietabelle:
$$
\begin{array}{cc|c|c}
 & x< \ln{2} & x=\ln{2} &  x> \ln{2} \\
\hline
(e^x-2)^2 & + & 0 & + \\
I'(x) & + & 0 & + \\
\hline
I(x) & \nearrow & \rightarrow & \nearrow \\
 & sms & Terrassenpunkt & sms 
\end{array}
$$
Die Funktion $I(x)$ ist auf ganz $\mathbb{R}$ streng monoton steigend und besitzt einen Terrassenpunkt bei $T(\ln2 | 0 )$.

\noindent {\bf  Aufgabe 2. b)}

$$ F'(x)= 0,5 \cdot 2 \cdot e^{2x} -4e^x+4 = (e^x-2)^2$$

\noindent {\bf  Aufgabe 2. c)}
Bestimmung der oberen Integrationsgrenze:
$$ 
\begin{array}{rcl}
 (e^x-2)^2 &=& 4 \\
e^x-2 &=&\pm 2 \quad \textnormal{ nur positive Lösung möglich}\\
e^x&=&4 \\
x &=& \ln{4}
\end{array}
$$
$$
\begin{array}{rcl}
A(u)&=&\int \limits_u^{\ln{4}}(4- (e^x-2)^2 )~{\rm d} x = \left[ 4x-0,5\cdot e^{2x}+4e^x-4x \right]_u^{\ln{4}}\\
A(u)&=&-0,5\cdot 16 + 4 \cdot 4 -0,5 \cdot e^{2u} - 4e^{u}\\
&=& 8-0,5 \cdot e^{2u} - 4e^{u}
\end{array}
$$


\noindent {\bf  Aufgabe 2. d)}
$$\lim \limits_{u \to -\infty}  8-0,5 \cdot e^{2u} - 4e^{u}=8-0-0=8$$
Die vom Graphen und der Geraden $y=4$ eingeschlossene Fläche hat einen endlichen Flächeninhalt und beträgt $8~FE$.

\noindent {\bf  Aufgabe 3.}\\
$I(x)$ wird durch Graph {\bf B} dargestellt. Graph {\bf A} besitzt zwei Terrassenpunkte, diese würden zwei doppelten Nullstellen von $G_f$ entsprechen. Diese sind jedoch nicht vorhanden.
Graph {\bf C} hat 5 Extrema, ist also mindestens 6. Ordnung, in Frage kommt aber nur eine Funktion der Ordnung 3+1, also 4. 

\begin{center}
	\includegraphics[width=7 cm]{../../viecher/endcomic.pdf}
	
	Hier geht es zurück zum \href{https://www.okuyakl.de/math/m12intapL222/aa222.pdf}{Aufgabenblatt}
\end{center}

\end{document}